% =====================================================================
% Capa e abertura
% =====================================================================
\thispagestyle{empty}
\begin{center}
\vspace*{2.2cm}
{\color{corcerrado}\scshape\large Caderno de preparação}\\[3pt]
{\color{corregra}\rule{0.55\textwidth}{1pt}}\\[26pt]
{\color{cortinta}\Huge\bfseries A marcha ao norte}\\[14pt]
{\color{cortinta}\Large\itshape o que você precisa saber\\[4pt] para defender esta qualificação}\\[34pt]
{\color{corregra}\rule{0.35\textwidth}{0.6pt}}\\[22pt]
{\large Victor Alves Rodrigues Amaral}\\[6pt]
{\normalsize PPGCIAMB / UFG}\\[26pt]
\begin{minipage}{0.80\textwidth}
\small\itshape\color{cortinta}
Companhia do texto de qualificação \emph{A marcha ao norte: dinâmica de uso e
cobertura da terra em Goiás (1985--2024) e seus fatores socioeconômicos} ---
115 páginas, 8 figuras, 38 referências, 58 rotinas, 31 decisões metodológicas.
Este caderno não substitui aquele documento: destrincha a matemática que ele
resume, ordena os números que ele espalha e ensaia as perguntas que ele
convida.
\end{minipage}
\vfill
{\footnotesize\color{corcerrado}Goiânia · agosto de 2026}
\end{center}
\clearpage

% ---------------------------------------------------------------------
\thispagestyle{empty}
\vspace*{0.6cm}
{\color{cortinta}\Large\bfseries Como usar este caderno}

\vspace{-4pt}
{\color{corregra}\rule{\textwidth}{1pt}}

\vspace{6pt}
Uma qualificação não se ganha no que se apresenta, e sim no que se responde. A
apresentação você controla; a arguição, não. O que sustenta os quarenta minutos
seguintes à sua fala é uma coisa só: saber \emph{por que} cada número saiu
daquele jeito e o que ele deixa de dizer. Este caderno foi montado para essa
segunda metade.

Ele tem sete partes e três modos de leitura.

\vspace{4pt}
\textbf{Leitura de fundo (uma semana antes).} Partes I a IV, na ordem. A Parte
III é a mais longa e a mais importante, pois é onde cada método aparece com a
fórmula, a tradução em português e, sobretudo, a fronteira do que ele autoriza
afirmar. É o que dá confiança para responder sem hesitar.

\textbf{Leitura de ensaio (dois dias antes).} Partes V e VI. A Parte V é a régua
do que se pode e não se pode dizer; a Parte VI traz cinquenta perguntas de banca
com resposta escrita. Responda cada uma em voz alta \emph{antes} de ler a
resposta sugerida.

\textbf{Leitura de véspera (na hora).} Parte VII: o roteiro de apresentação e o
cartão de bolso, que cabe numa página e traz os números que precisam sair da
boca sem consulta.

\vspace{10pt}
{\color{cortinta}\bfseries As quatro caixas.} Elas se repetem por todo o
caderno, e cada uma tem um trabalho:

\vspace{4pt}
\begin{naodiz}[O que o método NÃO diz]
A fronteira do instrumento. Ler esta caixa é o que impede a resposta que concede
demais e a que promete demais.
\end{naodiz}

\begin{nabanca}[Como dizer isto na banca]
Uma formulação pronta, em voz de quem domina o assunto e conhece os limites
dele. Não decore; entenda o movimento da frase.
\end{nabanca}

\begin{ancora}[Números-âncora]
Os valores que sustentam aquela afirmação. São os que valem memorizar.
\end{ancora}

\begin{armadilha}[Armadilha]
O erro que a pergunta convida a cometer --- em geral um excesso de afirmação que
soaria bem e não se sustenta.
\end{armadilha}

\vspace{8pt}
Ao longo do texto, a marca \onde{§4.2} indica a seção do documento de
qualificação em que aquilo está escrito, para que você possa ir à fonte durante
a defesa se precisar.

\clearpage
\pagenumbering{roman}
{\color{cortinta}\Large\bfseries Sumário}
\vspace{-6pt}
\makeatletter\renewcommand{\tableofcontents}{\@starttoc{toc}}\makeatother
\tableofcontents
\clearpage
\pagenumbering{arabic}
